% ═════════════════════════════════════════════════════════════════════════════
\section{Parameterized Verification of Quantum Circuits}
% ═════════════════════════════════════════════════════════════════════════════

The preceding sections have developed parameterized verification in the
classical setting: processes communicate through shared memory or message
passing, the parameter~$n$ counts the number of processes, and correctness
is expressed as a safety or liveness property over an infinite family of
transition systems. A structurally analogous challenge arises in quantum
computing. Most useful quantum algorithms, Shor's factoring algorithm,
Grover's search, the Quantum Fourier Transform, are \emph{parameterized by
input size}: a quantum \emph{program} takes~$n$ as input and dynamically
generates the corresponding $n$-qubit circuit. Verifying such a program
requires establishing correctness for all values of~$n$ simultaneously, a
goal that parallels classical parameterized verification but requires
substantially different technical machinery.

\paragraph{The verification problem.}
The correctness criterion for quantum programs is characteristically different
from the safety and liveness properties of Sections~2--8. Rather than asking
whether a bad state is ever reached, one asks whether the program computes
the right function: given a set of quantum state pairs $(p, q)$, does the
program, for every input size~$n$, correctly map input state~$p$ to output
state~$q$? This \emph{relational} (or \emph{functional}) specification captures
input-output correctness across all parameterized instances simultaneously.
It also subsumes \emph{equivalence checking} between two parameterized programs
--- a property of direct practical importance for verifying that an optimized
circuit implementation is functionally identical to a reference specification,
regardless of~$n$.

Existing approaches to quantum program verification have made notable progress
but fall short of fully automatic treatment of general parameterized circuits.
The Qbricks project~\cite{Qbricks} develops a deductive framework for quantum
programs with rich specifications, but verification is not fully automated and
requires user-provided proof annotations. The AutoQ
tool~\cite{AutoQ1,AutoQ2} supports automated verification of quantum circuits
using tree automata but is limited to a restricted class of parameterized
instances. A fully automatic framework for general parameterized quantum
programs was an open problem until recently~\cite{AutoQPara}.

\paragraph{Representing infinite families of quantum states.}
The central challenge in all parameterized verification is finding a finite
symbolic representation of an infinite family of system states, one state
space for each value of~$n$. In the classical setting, regular languages over
a finite alphabet serve this role (Section~4): a single finite automaton
compactly represents all reachable configurations across all instance sizes.
In the quantum setting the objects are quantum states rather than process
configurations. An $n$-qubit quantum state is a vector of $2^n$ complex
amplitudes, and a parameterized algorithm generates an infinite family of
such vectors, one for each~$n$, with a structured recursive form that grows
exponentially in size but must be captured in a finite representation.

The framework of~\cite{AutoQPara} achieves this through a new automata model
called \emph{synchronized weighted tree automata} (SWTAs). The starting point
is the observation that an $n$-qubit quantum state can be naturally represented
as a binary tree of depth~$n$: each internal node corresponds to a qubit, and
each leaf stores the complex amplitude associated with a particular
computational basis state. Standard tree automata recognize sets of such
trees, providing a symbolic representation of sets of quantum states that is
analogous to finite automata over words. SWTAs extend this with two additional
features. First, \emph{colors} attached to tree nodes serve as synchronization
labels, allowing the automaton to correlate corresponding positions across two
different state trees; this is the mechanism by which input-output pairs are
linked. Second, \emph{weights} on transitions encode complex amplitudes,
enabling the automaton to represent not merely which quantum states are
possible, but what amplitude each state carries. Together, these extensions
yield a model that defines a \emph{function} from color sequences to quantum
states, rather than simply a set of states. This functional view is essential
for relational verification: it allows a specification to assert what output
state is expected for each input state, across all input sizes simultaneously.

\paragraph{Transducers for quantum gate semantics.}
In regular model checking, the transition relation of a parameterized system
is encoded as a finite-state transducer, and reachability reduces to computing
the transitive closure of that transducer (Section~4). The quantum analogue
replaces the transition relation with a \emph{quantum gate}, and the transducer
must encode the linear transformation that the gate applies to the quantum
state tree. The framework of~\cite{AutoQPara} introduces a class of transducers
that operate directly on SWTAs: given an input SWTA representing a family of
quantum states, applying a gate transducer produces an output SWTA representing
the family of states after the gate has been applied.

A central concern in any symbolic framework for quantum computation is
\emph{succinctness}: a quantum gate on $n$ qubits is a $2^n \times 2^n$
unitary matrix, and a naive representation is exponential in~$n$. The
transducer representation avoids this blow-up. As an illustration, the
transducer for the Quantum Fourier Transform on $n$ qubits has a number of
states that is \emph{polynomial} in~$n$, exploiting the recursive structure
of the QFT to obtain an exponentially more compact representation than a
matrix encoding would require. This mirrors the advantage of transducers
over explicit transition tables in classical regular model checking.

A further key ingredient is a \emph{composition operator} for transducers.
Given transducers for two sub-circuits, the operator automatically constructs
a transducer for their sequential composition. Starting from transducers for
individual gates, repeated composition assembles the transducer for an
entire circuit without user intervention. The operator extends to
\emph{parameterized programs}: given a transducer for a finite circuit
constituting one layer of a parameterized algorithm, composition generates
a transducer for the full parameterized version with an arbitrary number
of such layers. This step, lifting from finite-circuit verification to
parameterized verification, is the direct quantum analogue of computing
the transitive closure of a transducer in classical regular model checking.

\paragraph{Verification algorithms and experimental results.}
With SWTAs as the specification formalism and transducers as the program model,
functional verification reduces to checking \emph{SWTA inclusion}: given the
SWTA~$\mathcal{A}$ specifying the desired input-output relation and the
SWTA~$\mathcal{A}'$ obtained by applying the program transducer to the input
specification, verify that $\llbracket \mathcal{A}' \rrbracket \subseteq
\llbracket \mathcal{B} \rrbracket$. Equivalence checking between two
parameterized programs reduces to an equality check between the output SWTAs
they produce when applied to a complete set of linearly independent inputs;
correctness of this reduction follows from the fact that quantum programs
implement linear transformations, so agreement on a complete input basis
implies agreement on all inputs. Both inclusion and equality of SWTAs are
decidable, and the framework provides practical algorithms for both.

The implementation, called \textsc{AutoQ-Para}~\cite{AutoQPara}, was evaluated
on a diverse set of benchmarks spanning algorithmic, arithmetic,
error-correction, and simulation circuits. It successfully verified the
Bernstein-Vazirani algorithm, a ripple-carry adder, a syndrome extraction
circuit from a quantum error-correcting code, and optimized implementations
of Grover's search and Hamiltonian simulation, comparing them against their
unoptimized reference counterparts. Verification times ranged from
14 milliseconds to 11 seconds, and to the best of the authors' knowledge
these are the first fully automatic parameterized verifications of circuits
of this kind.

\paragraph{Connections to classical parameterized verification.}
The structural parallels with classical parameterized verification are
instructive. In both settings the parameter~$n$ indexes an infinite family of
systems, and the goal is to verify a property of the entire family from a
finite symbolic representation. Regular languages and SWTAs both fill this
role compactly, words for classical configurations and trees for quantum states.
Transducers encode transition relations in the classical case and gate semantics
in the quantum case, and computing transitive closures or repeated compositions
is the operation that lifts finite reasoning to the parameterized family in
both. The SWTA inclusion check mirrors the backward reachability pattern of
monotonic abstraction, amounting to verifying that no reachable output state
lies outside the specification. What distinguishes the quantum setting is the
algebraic structure of the state space: quantum states carry complex amplitudes
subject to normalization and unitarity constraints, and gates are unitary linear
transformations rather than arbitrary relations. The SWTA model is designed to
respect this structure, with weights on transitions as the mechanism for
tracking amplitudes within an automata-theoretic framework. The result is, in
effect, a quantum lifting of regular model checking, adapted to the
linear-algebraic semantics of quantum computation.
